\section{Introduction}\label{sec:intro}

Central banks, finance ministries and firms in commodity-dependent economies use inflation forecasts, and a forecast is useful for decisions only together with a statement of its uncertainty. A prediction interval with a nominal level of 90 percent makes a specific claim: in about nine cases out of ten the realised value should fall inside it. Whether this claim holds in practice is the question of calibration, and it can be checked directly against realised data. It is not guaranteed by the forecasting model, because the model's own error distribution is estimated from a past that may not resemble the future.

Commodity-dependent economies are a demanding setting for calibration. Their price levels are exposed to exchange-rate movements and to commodity prices, and their inflation regimes can change within a few years. We do not cite systematic evidence for this statement; it is the motivation for the setting and not a result of the paper, and the paper does not compare commodity-dependent economies with other economies (Section~\ref{sec:scope}). A calibration made once on a past window may fail later, and a method that adapts the interval width online is a natural candidate. Conformal prediction \citep{vovk2005algorithmic} gives distribution-free finite-sample coverage under exchangeability, which time series violate. Adaptive conformal methods \citep{gibbs2021aci,gibbs2024faci,angelopoulos2023pid,bhatnagar2023saocp} were designed for that case, and a recent regime-aware variant, DRACP \citep{oancea2026dracp}, reports gains on euro-area inflation. Whether these results transfer to commodity-dependent economies, and at multi-month horizons, has not been tested to our knowledge; no systematic literature search was run, and this statement rests on the author's reading only.

\subsection{Problem and gap}

The problem addressed is the calibration of prediction intervals for the cumulative log growth of the consumer price index at horizons of 1, 3, 6 and 12 months in economies whose inflation is exposed to commodity and exchange-rate shocks. The gap is an evidence gap, not a methodological one. The closest work, \citet{oancea2026dracp}, evaluates DRACP on 28 harmonised euro-area and EU inflation series, 13 US macroeconomic and energy series and 7 daily financial series, mainly one step ahead, and states that its advantage is specific to one-step-ahead forecasting and depends on the point forecaster. What is missing is a benchmark of the same family of methods on a different panel, at multi-month horizons with outcomes that arrive with a delay, with decision rules fixed before the test run so that null outcomes can be reported as such. We provide that benchmark. We do not propose a new interval method, and we do not claim novelty beyond a controlled transfer of published methods to a new setting and a benchmark design that reports power and null results explicitly.

\subsection{What is done}

We construct a rolling-origin benchmark for the cumulative $h$-month log growth of the consumer price index at horizons 1, 3, 6 and 12 months. Six interval methods (split, rolling, ACI, FACI, conformal PID, DRACP) are placed on top of point forecasts, with a nominal level of 0.90 and with outcome feedback delayed by $h$ months, as it is in real use. The panels comprise 13 economies used to design the study and 7 hold-out economies for which no forecast had been computed before the pre-specified hold-out run. The primary outcome is the calibration error, defined per country as the absolute difference between empirical coverage and 0.90 and averaged over countries. Pooled coverage is reported but not used as the primary outcome, because opposite deviations of different countries can cancel in a pooled average. The decision rules for the hold-out run (R1 to R5) were written in the repository before the run, and the outcomes are reported exactly as the rules classify them. ``Pre-specified'' here refers to a self-attested commit history, not to an external registry, and the rules were written after the exploratory results were known (Section~\ref{sec:prespec}). An extension, specified in the eighth protocol entry before any forecast was computed for its countries, adds a second hold-out panel of 10 economies, a zero-shot foundation-model forecaster (Chronos) as a fourth base model, and the rules E1--E5.

\subsection{Research questions}

The study asks four questions.
\begin{enumerate}
\item[\textbf{RQ1}] Are adaptive conformal intervals (ACI, FACI, conformal PID, DRACP) better calibrated than a conformal calibration performed once (split) and than the native intervals of the forecasting model, in commodity-dependent economies and at horizons of 1 to 12 months?
\item[\textbf{RQ2}] Do adaptive methods improve on the simpler alternative of recomputing the conformal quantile from a rolling window of recent residuals (rule R2), and does a regime-aware method (DRACP) improve on ACI or is it equivalent to it (rule R3)?
\item[\textbf{RQ3}] How much of the under-coverage at the 12-month horizon is caused by the delay with which outcomes become known (rule R4)?
\item[\textbf{RQ4}] Do the conclusions depend on the panel of countries, on the base forecaster, on the warm-up size, and on the choice of a zero-shot foundation model as the point forecaster?
\end{enumerate}

\subsection{Summary of findings, stated plainly}

\begin{enumerate}
\item Adaptive methods were better calibrated than split conformal calibrated once on the first months and never updated, and better calibrated than the native intervals of the forecasting model (secondary, post hoc). The direction is the same in the exploratory data, on the first hold-out panel at both warm-up sizes, for three base forecasters, and after removing any single country. The registered rule R1, however, was met only at the 96-month warm-up, which is a sensitivity run and is not promoted to the primary result. At the primary 48-month warm-up the result for ACI minus split was $-3.0$~pp with $p=0.054$ at $h=3$ and $-2.8$~pp with $p=0.092$ at $h=6$ (Holm-adjusted 0.108 at both horizons); the confidence intervals, \ci{-5.4}{+0.5} and \ci{-3.9}{+1.2}, include zero.
\item An improvement of ACI over a rolling-window conformal method was not shown (R2) on the design panel and on the first hold-out panel (minimum detectable difference undefined, inconclusive). Point estimates are negative (about 1 to 3~pp of calibration error in most settings) and their intervals include zero in the primary run. On the second hold-out panel the same rule was met at warm-up 48 in part of the cells, and on the earlier 20 countries re-run from 2010-01 it is not significant, with the minimum detectable difference undefined in every cell, so inconclusive. The conclusion depends on the panel; the reason is not known and was not tested, and the panels are not averaged into one claim. No efficiency advantage in interval width or interval score is established.
\item For DRACP against ACI the outcome is inconclusive (R3): equivalence within 2~pp was shown at warm-up 48 and not at warm-up 96. No advantage of DRACP over ACI was detected in this configuration. Because the 2~pp margin is about the size of ACI's own calibration error, this verdict is weak by construction.
\item All four primary contrasts of the first hold-out panel are inconclusive under the power rule R5, because the minimum detectable difference is undefined on the full hold-out sample.
\item At the 12-month horizon pooled coverage is below nominal for every method. Delayed feedback accounts for part of the shortfall, but the pre-specified threshold (R4) was not met for ACI: $+2.4$~pp against 3~pp required, with an interval that includes zero.
\item On the second hold-out panel at warm-up 48, R1 was met for all four base models, so the advantage over a calibration performed once replicated on new countries at the primary warm-up; at warm-up 96 it was not met, on 41 and 35 pooled months. DRACP had a lower CE than ACI in every cell by 0.5 to 1.8~pp, none of the differences reached Holm $p$ below 0.05, and R3 stays inconclusive.
\item Chronos used zero-shot as a base forecaster did not change the calibration conclusions. On target months from 2025-01 its mean absolute error was lower than that of the naive baseline at every horizon, and no paired interval favours LightGBM or ARIMA over Chronos; this is descriptive, rests on 22 calendar months, and depends on an assumed, unverified pretraining cutoff. The language-model explainer has not been evaluated.
\end{enumerate}

\subsection{What is not claimed}

This work does not show that adaptive methods are better than rolling windows, it does not show that they are not (the answer depends on the panel), and it does not show that DRACP fails or does not transfer. Most pre-specified rules returned null or inconclusive outcomes, and a null outcome here is not evidence of equality. A contrast whose minimum detectable difference is undefined, or above 7~pp, is called inconclusive and never ``no effect''. The hold-out countries are untouched with respect to countries only (Section~\ref{sec:boundary}).

\subsection{Contributions}

\begin{enumerate}
\item A calibration benchmark on 30 economies in three panels, with calibration error, interval score, width and power-aware inference, where the null outcomes were specified in advance.
\item Evidence that adaptive methods are better calibrated than a calibration performed once and than the native intervals of the forecasting model, with a direction consistent across samples; the registered rule R1 was met on the first hold-out panel only at the sensitivity warm-up.
\item A secondary, post hoc comparison with native ARIMA and ETS intervals, with no multiplicity correction: ACI with ARIMA had a lower CE than the native ARIMA intervals at every horizon in both samples, and on the first hold-out panel the interval for the difference includes zero at $h=6$ and $h=12$.
\item A diagnostic of the effect of delayed feedback at the 12-month horizon.
\item An extension on a second hold-out panel of 10 economies with four base models including a zero-shot foundation model, specified before any forecast was computed for these countries; it replicates the advantage over one-shot calibration at the primary warm-up and shows that the conclusion about ACI against a rolling window depends on the panel, a result that is reported next to the earlier ones and not resolved.
\item A descriptive comparison of the point accuracy of Chronos with LightGBM, ARIMA and naive B on target months it cannot be assumed to have seen.
\item An open pipeline with dated hashed data snapshots, tests and a committed protocol. A platform module that serves inflation intervals and an explanation layer are engineering artefacts and not scientific contributions (Section~\ref{sec:architecture}).
\end{enumerate}

The remainder of the paper is organised as follows. Section~\ref{sec:related} reviews related work. Section~\ref{sec:method} formalises the problem and the methods. Section~\ref{sec:architecture} describes the pipeline and, briefly, the platform layers. Section~\ref{sec:data} describes the data and the experimental design, including what was fixed before which run. Section~\ref{sec:results} gives the results, Section~\ref{sec:discussion} discusses them, Section~\ref{sec:limitations} lists limitations, and Section~\ref{sec:repro} describes how to reproduce the work.
